\documentclass[11pt]{scrartcl}

\usepackage[stixtwo]{fontsetup}
\usepackage{amsmath}

\begin{document}
\[
	a_n = \begin{cases}
		\left(
		\dfrac{\left(1 - \frac{9}{n}\right)^n}{\left(1 + \frac{4}{n}\right)^n}
		\right)^3 \to e^{-39},
		 & \text{ha } n \text{ páros}     \\[3ex]
		-\left(
		\dfrac{\left(1 - \frac{9}{n}\right)^n}{\left(1 + \frac{4}{n}\right)^n}
		\right)^3 \to -e^{-39},
		 & \text{ha } n \text{ páratlan}
	\end{cases}
\]
A sorozat torlódási pontjai a $\pm e^{-39}$, vagyis
$\limsup a_n = e^{-39} \ne -e^{-39} = \liminf a_n$, és a sorozat nem
konvergens.

Mivel
\(
\left| \frac{9 - n}{4 + n} \right|^n \to e^{-13}
\),
ezért elég nagy $n$ esetén
\[
	|b_n|
	= \left( \left| \frac{9 - n}{4 + n} \right|^n \right)^{n^2}
	< \left( e^{-13} + \varepsilon \right)^{n^2}
	\to 0,
\]
ha $0 < \varepsilon < 1 - e^{-13}$. Így a sorozat egyetlen torlódási pontja
a $0$, és $\limsup b_n = \liminf b_n = 0 = \lim b_n$.
\end{document}
